\documentclass[11pt,twoside]{amsart}
\usepackage{amssymb, amsmath, enumerate, palatino, hyperref}
\usepackage{caption}
\usepackage{fullpage}

\usepackage[T1]{fontenc}
\renewcommand{\labelitemi}{\guillemotright}
\usepackage{mathrsfs}
\hypersetup{hidelinks}

\theoremstyle{plain}
\newtheorem{prop}{Proposition}
\newtheorem{lemma}[prop]{Lemma}
\newtheorem{thm}[prop]{Theorem}
\newtheorem{cor}[prop]{Corollary}
\theoremstyle{remark}
\newtheorem{rmk}[prop]{Remark}
\newtheorem*{rmk*}{Remark}
\newtheorem{prob}{Problem}
\newtheorem*{bonus*}{Bonus Problem}
\theoremstyle{definition}
\newtheorem{ex}[prop]{Example}
\newtheorem{defn}[prop]{Definition}

\newcommand{\cA}{\mathcal A}
\newcommand{\cZ}{\mathcal Z}
\newcommand{\SEQ}{\operatorname{SEQ}}
\newcommand{\MSET}{\operatorname{MSET}}
\newcommand{\PSET}{\operatorname{PSET}}
\newcommand{\FF}{\mathbb F}
\newcommand{\qbinom}[2]{\genfrac{[}{]}{0pt}{}{#1}{#2}_{\!q}}

\title{Math 372: Combinatorics\\ Homework 04}
\author{Due Friday, 30 October, 10\textsc{p.m.}, via Gradescope}

\begin{document}
\maketitle

\begin{rmk*}
\textbf{When you can start.} Problem 1 uses the constructions from Monday, October~12, and
Problem~2 uses the multisets and partitions from Wednesday, October~14: do these two before or
over fall break. Problem~3 needs Monday, October~26 (Young diagrams and $q$-binomial
coefficients, which is also the break reading), and Problem~4 is best done after Wednesday,
October~28. The second half is deliberately short.
\end{rmk*}

\begin{rmk*}
\textbf{Conventions.} A \emph{composition} of $n$ is a finite sequence of positive integers
with sum $n$; a \emph{partition} of $n$ is a multiset of positive integers with sum $n$, written
as a weakly decreasing sequence $\lambda=(\lambda_1\ge\lambda_2\ge\cdots)$. The empty
composition and the empty partition are the unique objects of size~$0$. Fibonacci numbers are
$F_0=0$, $F_1=1$, $F_{n+1}=F_n+F_{n-1}$. For $0\le k\le N$, let $L(k,N-k)$ be the set of
partitions with at most $k$ parts, each at most $N-k$ (the Young diagrams that fit in a box with
$k$ rows and $N-k$ columns), ordered by containment of diagrams, and define
\[
\qbinom Nk=\sum_{\lambda\in L(k,N-k)}q^{|\lambda|},\qquad |\lambda|=\lambda_1+\lambda_2+\cdots.
\]
You may use the recurrence from class on October~26,
$\qbinom Nk=\qbinom{N-1}{k}+q^{N-k}\qbinom{N-1}{k-1}$ for $0<k<N$, together with
$\qbinom N0=\qbinom NN=1$.
\end{rmk*}

\begin{rmk*}
\textbf{On write-ups.} Each problem is worth up to five points for mathematical content and up to
two for writing. Write for another student in this course: say what you are going to show before
you show it, justify each step, and use complete sentences. Whenever you give a specification,
say in words what the objects are before writing the generating function. End your submission
with a short \emph{Methods and tools} note: what software you used and for what, and how you
checked anything a machine produced for you.
\end{rmk*}

%%% Problem 1 %%%
\begin{prob}[specifications]
\begin{enumerate}[(a)]
\item Let $\cA$ be the class of binary words with no three consecutive $1$s, size $=$ length.
Give a specification of $\cA$ in terms of $\SEQ$, $+$ and $\times$, deduce its generating
function $A(z)$ as a single fraction, and compute $a_0,\dots,a_6$ from it. Check $a_4$ by
listing the words of length $4$ that are \emph{not} in $\cA$.
\item Let $\mathcal C$ be the class of compositions in which every part is at least $2$.
Give a specification, find $C(z)$, and show that for $n\ge2$ the number $c_n$ equals the number
of compositions of $n-2$ into parts $1$ and $2$. Then find a bijection that explains this.
(Suggestion: rewrite each part $a\ge2$ of a composition in $\mathcal C$ as $2$ followed by
$a-2$ ones.)
\item Mark each part of a composition with a variable $u$: show that
\[
\sum_{n,k\ge0}\#\{\text{compositions of $n$ with exactly $k$ parts}\}\,u^kz^n
=\frac{1}{1-\dfrac{uz}{1-z}},
\]
and deduce that the number of compositions of $n\ge1$ with exactly $k$ parts is
$\binom{n-1}{k-1}$. Explain the same count with the ``dots and bars'' picture from class, and
check that summing over $k$ gives the total you already know.
\end{enumerate}
\end{prob}

%%% Problem 2 %%%
\begin{prob}[partitions]
\begin{enumerate}[(a)]
\item Prove, using generating functions, that for every $n$ the number of partitions of $n$
with no part divisible by $3$ equals the number of partitions of $n$ in which no part appears
three or more times.
\item Check (a) for $n=6$ by listing both sets of partitions.
\item State and prove the analogous statement with $3$ replaced by any integer $d\ge2$. What
statement from class is the case $d=2$?
\end{enumerate}
\end{prob}

%%% Problem 3 %%%
\begin{prob}[$q$-binomial coefficients]
\begin{enumerate}[(a)]
\item Prove the product formula
\[
\qbinom Nk=\frac{(1-q^N)(1-q^{N-1})\cdots(1-q^{N-k+1})}{(1-q^k)(1-q^{k-1})\cdots(1-q)},
\]
by showing that the right-hand side satisfies the recurrence and boundary values in the
Conventions. (In particular the right-hand side is a polynomial, which is not obvious.)
\item Prove the second recurrence $\qbinom Nk=q^{k}\qbinom{N-1}{k}+\qbinom{N-1}{k-1}$ by
sorting the diagrams in $L(k,N-k)$ according to their first column.
\item Show that the coefficients of $\qbinom Nk$ are symmetric:
$q^{k(N-k)}\qbinom Nk\big|_{q\to q^{-1}}=\qbinom Nk$. (Find an involution of $L(k,N-k)$.)
\end{enumerate}
\end{prob}

%%% Problem 4 %%%
\begin{prob}[unimodality with no linear algebra, for two rows]
Fix $n\ge1$ and consider the poset $L(2,n)$: partitions $(\lambda_1,\lambda_2)$ with
$n\ge\lambda_1\ge\lambda_2\ge0$, graded by $\lambda_1+\lambda_2$, of rank $2n$.
\begin{enumerate}[(a)]
\item Show that the number of elements of rank $i$ is $\lfloor i/2\rfloor+1$ for $0\le i\le n$,
and use (3c) or a direct argument for $n\le i\le2n$. Conclude that $L(2,n)$ is rank-symmetric
and rank-unimodal, and write down its rank sizes when $n=4$.
\item Find order-matchings $L(2,n)_0\to L(2,n)_1\to\cdots\to L(2,n)_n$ and
$L(2,n)_n\leftarrow L(2,n)_{n+1}\leftarrow\cdots\leftarrow L(2,n)_{2n}$, and verify that they
are order-matchings. (Change one part at a time.)
\item Conclude that the largest antichain in $L(2,n)$ has $\lfloor n/2\rfloor+1$ elements, and
draw the chains your matchings produce when $n=4$.
\end{enumerate}
\end{prob}

%%% Bonus %%%
\begin{bonus*}[subspaces]
Let $q$ be a prime power and $\FF_q$ the field with $q$ elements. Show that the number of
$k$-dimensional subspaces of $\FF_q^N$ is $\qbinom Nk$ evaluated at that $q$. (Count ordered
bases two ways.) How many $2$-dimensional subspaces does $\FF_2^4$ have?
\end{bonus*}

\end{document}
